\section{Electronics and Sensing}
\label{sec:electronics}

This section is short by intent. It records what the design notes, the photographs of the
robot and the simulation model establish about the electronics, and nothing beyond that. No
circuit diagram or parts list of the competition robot was available when this report was
written.

\subsection{Actuators}
Two N20 gearmotors with encoder boards drive the wheels, and one MG995 servo drives the
gripper (\cref{sec:drive,sec:gripper}). The motors were selected for their size.

\subsection{Sensors}
\begin{description}
  \item[Infrared array.] Eight analog reflectance channels at a pitch of 15\unit{mm}, read as
        voltages between 0\unit{V} on a black surface and 5\unit{V} on a white one. The
        array serves the line following and, in the simulated mission, the final approach to
        every station.
  \item[Depth camera.] An Intel RealSense D435, which provides a colour image and a depth
        image. In the simulation it is the only sensor used for mapping, localisation,
        obstacle detection and the reading of QR codes.
  \item[Wheel encoders.] One on each motor, the source of wheel odometry.
  \item[Inertial sensor.] The simulation model carries an inertial sensor of the MPU6050
        class, flat on the plate between the wheels, whose yaw rate is fused with the wheel
        odometry (\cref{sec:odometry}).
\end{description}

\subsection{Power and Control}
The photograph of the robot (\cref{fig:v2-intro}b) shows a battery with an XT60 connector
at the rear of the plate, a step-down converter module, and a
microcontroller board on a prototyping board with a solderless breadboard for the signal
wiring. Both the electronics and the mechanical design were checked extensively before the
competition.

\subsection{The Failure at the State Round}
In the state round the infrared array stopped working and the robot could not follow the
line, so the task was not completed. The cause was not
established afterwards, and I do not speculate on it here. Two consequences are drawn for
the design. The line following depended on a single sensor with no fallback. And the photograph shows
signal wiring on a solderless breadboard, which is not a reliable connection on a moving
robot. A soldered or connectorised harness, and a check of all eight channels against known
black and white references at power-up, are the corrections proposed
in \cref{sec:discussion}.
